%! TeX root = ../thesis.tex
\chapter{Conclusions and outlook}

We have shown that by splitting the AIM into two components,
one part can be represented as bits,
and another can be represented as a vector of CI excitations.
This allows us to speed up the calculation $\ham \ket{\psi}$
of the Anderson Hamiltonian on a state by
up to a factor of 2000 depending on the number of bath sites and excitation
with an accuracy loss of $O(10^{-15})$ for the wave function.
Calculations which took tens of seconds can now be done in
a matter of hundreds of \unit{\milli\second}.

Using this new tool,
our goal is to determine properties and Green's functions of systems
which contain hundreds of bath sites which is going to be
the work of my master thesis.

Although the current implementation shows numerical consistence with the
existing bit representation,
some restrictions remain.
One is that hopping is currently restricted to real values,
which is not hard to extend to the complex domain if necessary.
Another factor is the addition of a spin dependency
to the mixed operator $\ham_{\mathrm{mix}}$.
Instead of using tuples $(i, j, t)$,
we can use $(i, j, \sigma, t)$ and adjust the backend which is not much work.
If this feature is realized,
all three terms of~\cref{eq:hamiltonian-split}
can accommodate spin-dependent energies and hopping.

We can extend the feature set to triple excitations if doubles are not
accurate enough. For fast unranking (see~\cref{eq:unranking}) this
requires solving a cubic equation on integers and storage of a matrix
which would take a few \unit{\giga\byte} of space (\cref{tab:ham-storage}).

At the moment, terms of opposite spin hopping, e.g., $\cdag_{i\up}\cc_{j\dn}$
will evaluate with a possible wrong sign in the bit component.
These terms are not necessary as long as we have time reversal symmetry,
since we can change into the $\hat{J}$, $\hat{J_z}$ basis.
If one wants to implement a few of this kind of terms,
it is possible to write a custom operator,
otherwise it would probably be better to change the whole underlying order from
spin-sorting defined in \cref{eq:sort-by-spin} to site-sorting
(\cref{eq:sort-by-site}),
which requires a major rewrite.

To get even more performance,
we could employ multithreading.
The first choice would be the application
of the vector Hamiltonian on the vector component $\ham_v \ket{v}$,
since it does not change dictionary keys
and each matrix-vector multiplication is independent of each other.
